\section{讨论}
\label{sec:discussion}

本文刻画了在传播强度随时期变化、后代数过度离散且仅有部分感染被报告的条件下，基于病例数的再生数估计误差如何随病例量变化。对当期值而言，误差随报告病例的累积无限下降。对跨期均值和一步预测而言，每期误差趋于环境方差；平衡规模 $\mx=A/v_R$ 标出了这样一个病例量，超过它之后病例数每加倍一次，至多只能消除剩余方差的四分之一。在四个流行季的每周流感住院数据中，两种表观统计量的截距均为正；经由数据匹配模拟的两步反演，两者都对应约 0.04 的潜伏环境方差。全国周分量约占截距的 40\%--50\%，而强制报告下报告完整度对截距的贡献小于截距 17\% 的等效界（按参考报告波动 0.10 设定）。

上述理论把种群生态学中的一个标准结果推广到监测数据。生态学模型把种群增长的方差分为两项：随种群规模下降的种群统计项，以及不随种群规模下降的环境项\citep{engen1998demographic, lande2003stochastic}。本文模型保留这一分解，并加入几个对监测有实质意义的要素。传播被设定为后代服从负二项分布的分支过程，超级传播因此放大种群统计项\citep{lloydsmith2005superspreading}。二项报告引入另一项随病例量下降的方差，报告波动则引入一项不随病例量下降的方差（命题~\ref{prop:report-noise}）。由于关注的量是跨期平均的估计，平衡规模还决定预算在每期病例数与时期数之间的最优划分（命题~\ref{prop:allocation}）。这一规则与两阶段抽样中的最优分配形式相同\citep{cochran1977sampling}，只是以平衡规模代替单元内方差与单元间方差之比。

本文结果与 Parag 等人的信息论分析相互补充\citep{parag2021improved, parag2022quantifying}。他们针对含噪声流行曲线提出的度量已经包含随时间变化的报告比例，后者通过各期报告概率的几何平均进入度量；其补充材料还讨论了负二项扩展\citep{parag2022quantifying}。他们的估计目标是逐期取值的序列；在泊松更新模型中，每期估计的误差方差为 $\Rbar/m$，报告的作用相当于对有效病例数重新标度。本文的估计目标是跨期均值，同一时期内所有传播链共有的变异会加入一项不随 $m$ 下降的方差。两种描述在病例数较少时一致，在平衡规模附近开始分离。在 $R_A$ 的表观尺度上，拟合方差在每周 22 例住院时为泊松方差的两倍，在 261 例时为十倍；在潜伏尺度上，相应位置为 28 例和 262 例。两种描述因此各有用途。泊松分析按各数据流所含当前传播信息的多少对其排序，本文的方差分解则指出继续增加病例何时不再改善对典型传播水平的估计。

两步反演把表观统计量与潜伏模型联系起来。模拟首先确认，两种统计量都不直接估计 $v_R/\Rbar^2$：$G$ 的截距是该量的 2.7 倍，$R_A$ 约为其一半，原因在于两者在一周之内汇总增长的方式不同。校准之后，两种统计量隐含的潜伏方差为 0.043 和 0.040（比值 0.93；95\% CI：0.78, 1.18）。若两者响应同一潜伏变异，且模拟正确刻画了各自的变换方式，就会出现这种一致。作为对校准的检验，这种一致的效力有限。在模拟的各设定中，只有以作用于报告计数的同步冲击代替传播冲击时一致性才被打破；在该设定下，数据中 $G$ 的滞后一阶自相关高于 99\% 的模拟面板。模拟面板还复现了数据的组内相关系数、全国周份额和滞后一阶自相关。在所述不确定性范围内，校准因此与经验截距相容，但需满足三个条件。隐含方差取决于真实代际间隔。在流感序列间隔 2.9--4.3 天的 95\% 区间内\citep{cowling2009estimation}，它从 0.026 变化到 0.063，反演识别的是以代际间隔为条件的 $v_R$。周冲击被设定为相互独立。观测到的自相关位于独立冲击下模拟分布的中部，对 $R_A$ 而言低于一阶系数为 0.25 时 97\% 的模拟面板；若一阶系数为 0.5，隐含方差会降至约 0.026。最后，校准后的理论 $b/a$ 低于两个实证点估计，后者位于模拟抽样分布的第 91 和第 93 百分位。周住院数的过度离散可以解释这一超出。方差均值比为 1.5 时，模拟的 $b/a$ 对 $R_A$ 升至 68.0、对 $G$ 升至 40.2，均高于实证值，而隐含方差保持在 0.040 和 0.044。真实代际间隔为 2.9 天时同样能复现实证 $b/a$，对应的隐含方差为 0.026。因此，这一超出可能来自方差均值比小于 1.5 的过度离散，也可能来自较短的代际间隔。两种解释对应的隐含方差分别约为 0.040 和 0.026。数据的自相关不支持较短的间隔，但不能将其排除。数据不能排除传播冲击与报告冲击混合的情形（命题~\ref{prop:report-noise}）。

实证结果最适合解读为单期病例量的收益递减尺度。以 $R_A$ 尺度上每周 37.4 例住院的表观平衡规模计，报告量已达平衡规模的州可获得病例量不受限时可达精度的一半，要达到 90\% 则需九倍于此的住院数。分析窗口中半数以上的州周已经位于这一尺度之上。这一极限针对的是单期病例量。对 $T$ 个独立时期合并的估计，其误差按 $v_R/T$ 下降；向先验均值收缩的估计量可以低于单期值（附录~\ref{app:shrinkage}）；在观测到的规模范围内，无截距幂律与仿射形式的拟合也大致相当。平衡规模因此回答的是收益从何处开始递减，并不意味着增加任何一类数据都不再降低误差。这一解读对截距的来源同样稳健。无论截距中传播、检测与统计量构造各占多少，只要某个误差分量在同一时期内共有且不随病例量下降，它就会限制每期增加病例的价值，并提高增加时期数的价值。

多尺度分解指出了这一分量所处的层级。$R_A$ 截距的约 41\% 与 $G$ 截距的 51\% 为同一日历周内各州所共有；对周重抽样使截距估计的变动幅度达到对州重抽样的 1.5--1.9 倍。全国周分量在同一周进入每个州的损失，单靠扩大某一州的监测规模无法缩小它。合并多州可以去掉州特有部分；在本数据中，合并 25--33 个州后，合并方差已接近全国周分量，约为单州方差的四分之一。这一份额与天气、学校校历和节假日等全国性影响相符，也与全国层面检测或收治实践的变化相符\citep{dalziel2018urbanization}。出于同样的原因，推断必须采用双向结构；仅对州重抽样所得的截距区间约窄了二至五倍\citep{cameron2011robust}。

在观察性比较所允许的范围内，两种报告制度的对比与较高机构参与度的益处相一致。强制报告下，报告医院占比的中位数约为 93\%；按报告完整度调整后，$R_A$ 的截距变化 2\%，上界为 5\%，在截距 17\% 的等效界内与零等效。2024 年自愿报告期报告占比中位数降至约 46\%，此时增长对数的四分位距是 2023 年同期的 1.58 倍。该比值的区间为 0.95--2.23，包含 1，且比较只依赖一个非流行季。几项检查针对的是主要混杂因素。2025 年对 2023 年的安慰剂比较给出 1.01，按报告规模和按日历月匹配后比值仍大于 1，同期 COVID-19 与 RSV 住院数据的比值也大于 1。这些结果与较低且多变的机构参与度伴随更不稳定的增长统计量相一致，但不能确认报告强制要求的因果效应。调整只度量了机构参与度，报告医院内部的检测强度、编码与就医行为并未纳入。实验室代理变量占 $R_A$ 截距的 23\%，但它无法区分传播与检出，因为实验室增长带有自身噪声，会使共有份额向零衰减。

这些结果为平均再生数的常规监测提出了一条设计规则。在各辖区每周住院数已超过平衡规模的情况下，把追加资源用于增加时期数和增加变异不同步的辖区（式~(\ref{eq:pooling})），比用于增加每期病例更能降低误差。完整报告的价值取决于系统实际存在的报告波动。在第~\ref{sec:res-theory}~节的假设示例中，相对波动 $\tau=0.15$，报告标准化对误差的降低幅度大于每期病例量增加十倍。强制报告下观测到的完整度上界只允许 $\tau$ 约为 0.04；取此值时，病例量增加十倍使均方误差由 $15.2\times10^{-3}$ 降至 $11.4\times10^{-3}$，报告标准化则降至 $14.6\times10^{-3}$。因此，在本数据中维持较高的机构参与度与较低的波动相伴，其相对于病例量的价值随参与不完整所带来的波动增大而上升。对于季间变异的判断，流行季是限制性单位。只有四个流行季时，季间方差的任何无偏估计量，其相对标准误都不低于 0.82（推论~\ref{cor:season}）。替代分母权重下的截距比未能区分机制，因为仅含报告噪声的流行增长已经会产生 $q_1<1<q_2$。

这些结论的适用范围由数据与校准决定。实证结果依赖一个监测系统、四个流感季和按规则界定的早期增长窗口。分析过程中窗口规则修订过一次，更新型比值改为主统计量；表~\ref{tab:s9} 列出了各设定及其角色与日期，主要估计值都落在彼此的区间之内。反演以代际间隔以及模拟的延迟与后代过程为条件；$R_A$ 是对周增长的变换，并非 $R_t$ 的后验估计\citep{cori2013new}。各命题假定环境变异在时期之间相互独立、原发病例数固定；第~\ref{sec:meth-loss}~节的更新方程表述表明，感染压力的不确定性会再增加一项不随病例量下降的方差。等效界使用各统计量的样本均值，规则窗口中截距比的零假设范围取自日历窗口模拟，图~\ref{fig:floor}B 的分箱误差线未考虑聚类。以下三个方向可以使结论更为确切：带显式入院延迟并采用外部估计代际间隔的更新估计；结合独立检测指标、含显式周效应的分层模型；以及另一个实行强制报告且覆盖更多流行季的监测系统。
